% Options for packages loaded elsewhere
\PassOptionsToPackage{unicode}{hyperref}
\PassOptionsToPackage{hyphens}{url}
%
\documentclass[
  11pt,
]{article}
\usepackage{amsmath,amssymb}
\usepackage{iftex}
\ifPDFTeX
  \usepackage[T1]{fontenc}
  \usepackage[utf8]{inputenc}
  \usepackage{textcomp} % provide euro and other symbols
\else % if luatex or xetex
  \usepackage{unicode-math} % this also loads fontspec
  \defaultfontfeatures{Scale=MatchLowercase}
  \defaultfontfeatures[\rmfamily]{Ligatures=TeX,Scale=1}
\fi
\usepackage{lmodern}
\ifPDFTeX\else
  % xetex/luatex font selection
\fi
% Use upquote if available, for straight quotes in verbatim environments
\IfFileExists{upquote.sty}{\usepackage{upquote}}{}
\IfFileExists{microtype.sty}{% use microtype if available
  \usepackage[]{microtype}
  \UseMicrotypeSet[protrusion]{basicmath} % disable protrusion for tt fonts
}{}
\makeatletter
\@ifundefined{KOMAClassName}{% if non-KOMA class
  \IfFileExists{parskip.sty}{%
    \usepackage{parskip}
  }{% else
    \setlength{\parindent}{0pt}
    \setlength{\parskip}{6pt plus 2pt minus 1pt}}
}{% if KOMA class
  \KOMAoptions{parskip=half}}
\makeatother
\usepackage{xcolor}
\usepackage[margin=1in]{geometry}
\usepackage{color}
\usepackage{fancyvrb}
\newcommand{\VerbBar}{|}
\newcommand{\VERB}{\Verb[commandchars=\\\{\}]}
\DefineVerbatimEnvironment{Highlighting}{Verbatim}{commandchars=\\\{\}}
% Add ',fontsize=\small' for more characters per line
\usepackage{framed}
\definecolor{shadecolor}{RGB}{248,248,248}
\newenvironment{Shaded}{\begin{snugshade}}{\end{snugshade}}
\newcommand{\AlertTok}[1]{\textcolor[rgb]{0.94,0.16,0.16}{#1}}
\newcommand{\AnnotationTok}[1]{\textcolor[rgb]{0.56,0.35,0.01}{\textbf{\textit{#1}}}}
\newcommand{\AttributeTok}[1]{\textcolor[rgb]{0.13,0.29,0.53}{#1}}
\newcommand{\BaseNTok}[1]{\textcolor[rgb]{0.00,0.00,0.81}{#1}}
\newcommand{\BuiltInTok}[1]{#1}
\newcommand{\CharTok}[1]{\textcolor[rgb]{0.31,0.60,0.02}{#1}}
\newcommand{\CommentTok}[1]{\textcolor[rgb]{0.56,0.35,0.01}{\textit{#1}}}
\newcommand{\CommentVarTok}[1]{\textcolor[rgb]{0.56,0.35,0.01}{\textbf{\textit{#1}}}}
\newcommand{\ConstantTok}[1]{\textcolor[rgb]{0.56,0.35,0.01}{#1}}
\newcommand{\ControlFlowTok}[1]{\textcolor[rgb]{0.13,0.29,0.53}{\textbf{#1}}}
\newcommand{\DataTypeTok}[1]{\textcolor[rgb]{0.13,0.29,0.53}{#1}}
\newcommand{\DecValTok}[1]{\textcolor[rgb]{0.00,0.00,0.81}{#1}}
\newcommand{\DocumentationTok}[1]{\textcolor[rgb]{0.56,0.35,0.01}{\textbf{\textit{#1}}}}
\newcommand{\ErrorTok}[1]{\textcolor[rgb]{0.64,0.00,0.00}{\textbf{#1}}}
\newcommand{\ExtensionTok}[1]{#1}
\newcommand{\FloatTok}[1]{\textcolor[rgb]{0.00,0.00,0.81}{#1}}
\newcommand{\FunctionTok}[1]{\textcolor[rgb]{0.13,0.29,0.53}{\textbf{#1}}}
\newcommand{\ImportTok}[1]{#1}
\newcommand{\InformationTok}[1]{\textcolor[rgb]{0.56,0.35,0.01}{\textbf{\textit{#1}}}}
\newcommand{\KeywordTok}[1]{\textcolor[rgb]{0.13,0.29,0.53}{\textbf{#1}}}
\newcommand{\NormalTok}[1]{#1}
\newcommand{\OperatorTok}[1]{\textcolor[rgb]{0.81,0.36,0.00}{\textbf{#1}}}
\newcommand{\OtherTok}[1]{\textcolor[rgb]{0.56,0.35,0.01}{#1}}
\newcommand{\PreprocessorTok}[1]{\textcolor[rgb]{0.56,0.35,0.01}{\textit{#1}}}
\newcommand{\RegionMarkerTok}[1]{#1}
\newcommand{\SpecialCharTok}[1]{\textcolor[rgb]{0.81,0.36,0.00}{\textbf{#1}}}
\newcommand{\SpecialStringTok}[1]{\textcolor[rgb]{0.31,0.60,0.02}{#1}}
\newcommand{\StringTok}[1]{\textcolor[rgb]{0.31,0.60,0.02}{#1}}
\newcommand{\VariableTok}[1]{\textcolor[rgb]{0.00,0.00,0.00}{#1}}
\newcommand{\VerbatimStringTok}[1]{\textcolor[rgb]{0.31,0.60,0.02}{#1}}
\newcommand{\WarningTok}[1]{\textcolor[rgb]{0.56,0.35,0.01}{\textbf{\textit{#1}}}}
\usepackage{longtable,booktabs,array}
\usepackage{calc} % for calculating minipage widths
% Correct order of tables after \paragraph or \subparagraph
\usepackage{etoolbox}
\makeatletter
\patchcmd\longtable{\par}{\if@noskipsec\mbox{}\fi\par}{}{}
\makeatother
% Allow footnotes in longtable head/foot
\IfFileExists{footnotehyper.sty}{\usepackage{footnotehyper}}{\usepackage{footnote}}
\makesavenoteenv{longtable}
\setlength{\emergencystretch}{3em} % prevent overfull lines
\providecommand{\tightlist}{%
  \setlength{\itemsep}{0pt}\setlength{\parskip}{0pt}}
\setcounter{secnumdepth}{5}
\ifLuaTeX
  \usepackage{selnolig}  % disable illegal ligatures
\fi
\usepackage{bookmark}
\IfFileExists{xurl.sty}{\usepackage{xurl}}{} % add URL line breaks if available
\urlstyle{same}
\hypersetup{
  hidelinks,
  pdfcreator={LaTeX via pandoc}}

\author{}
\date{}

\begin{document}

{
\setcounter{tocdepth}{3}
\tableofcontents
}
\section{Project 4 - Long Memory, ARMA, FIGARCH, and
Value-at-Risk}\label{project-4---long-memory-arma-figarch-and-value-at-risk}

\textbf{Course:} Financial Econometrics\\
\textbf{Author:} Omid Karami Chamgordani - 400203402\\
\textbf{Date:} 2023-02-07

\subsection{Objective}\label{objective}

This project studies daily returns of the financial institutions index
between 2008 and 2023. It examines long memory, selects an ARMA mean
equation, tests residual diagnostics and normality, estimates a
FI-GARCH-type volatility model, and computes Value-at-Risk using several
methods.

\subsection{Libraries and data}\label{libraries-and-data}

The cleaned Excel files contain daily dates and prices. The financial
institutions index is stored in \texttt{sarmaye.xlsx}; the overall
market index is stored in \texttt{index.xlsx} and is later used as an
external regressor.

\begin{Shaded}
\begin{Highlighting}[]
\FunctionTok{library}\NormalTok{(pdR)}
\FunctionTok{library}\NormalTok{(tidyverse)}
\FunctionTok{library}\NormalTok{(ggplot2)}
\FunctionTok{library}\NormalTok{(forecast)}
\FunctionTok{library}\NormalTok{(tseries)}
\FunctionTok{library}\NormalTok{(rio)}
\FunctionTok{library}\NormalTok{(readxl)}
\FunctionTok{library}\NormalTok{(zoo)}
\FunctionTok{library}\NormalTok{(lmtest)}
\FunctionTok{library}\NormalTok{(qqplotr)}
\FunctionTok{library}\NormalTok{(devtools)}
\FunctionTok{library}\NormalTok{(fracdiff)}
\FunctionTok{library}\NormalTok{(MASS)}
\FunctionTok{library}\NormalTok{(xts)}
\FunctionTok{library}\NormalTok{(PerformanceAnalytics)}
\FunctionTok{library}\NormalTok{(tidyr)}
\FunctionTok{library}\NormalTok{(rugarch)}
\FunctionTok{library}\NormalTok{(moments)}
\FunctionTok{library}\NormalTok{(dplyr)}
\FunctionTok{library}\NormalTok{(nortest)}
\FunctionTok{library}\NormalTok{(fBasics)}
\FunctionTok{library}\NormalTok{(TTR)}
\FunctionTok{library}\NormalTok{(quantmod)}
\FunctionTok{library}\NormalTok{(stargazer)}
\FunctionTok{library}\NormalTok{(ggthemes)}
\FunctionTok{library}\NormalTok{(gridExtra)}

\NormalTok{sarmaye }\OtherTok{\textless{}{-}} \FunctionTok{read\_excel}\NormalTok{(}\StringTok{"data/sarmaye.xlsx"}\NormalTok{)}
\NormalTok{ind }\OtherTok{\textless{}{-}} \FunctionTok{read\_excel}\NormalTok{(}\StringTok{"data/index.xlsx"}\NormalTok{)}
\end{Highlighting}
\end{Shaded}

\subsection{Price and return series}\label{price-and-return-series}

The price series is converted to a daily time series with frequency 235.
Returns are defined as first differences of log prices. In the plot,
daily returns remain within approximately -10 percent and +10 percent.

\begin{Shaded}
\begin{Highlighting}[]
\NormalTok{x }\OtherTok{\textless{}{-}} \FunctionTok{ts}\NormalTok{(sarmaye}\SpecialCharTok{$}\NormalTok{price, }\AttributeTok{start=}\FunctionTok{c}\NormalTok{(}\DecValTok{2008}\NormalTok{, }\DecValTok{12}\NormalTok{, }\DecValTok{06}\NormalTok{), }\AttributeTok{frequency=}\DecValTok{235}\NormalTok{)}
\FunctionTok{ts.plot}\NormalTok{(x, }\AttributeTok{main=}\StringTok{"daily price of investment index"}\NormalTok{, }\AttributeTok{ylab=}\StringTok{"price"}\NormalTok{)}

\NormalTok{returns }\OtherTok{\textless{}{-}} \FunctionTok{diff}\NormalTok{(}\FunctionTok{log}\NormalTok{(x), }\AttributeTok{lag=}\DecValTok{1}\NormalTok{)}
\NormalTok{returns2 }\OtherTok{\textless{}{-}} \FunctionTok{ts}\NormalTok{(returns, }\AttributeTok{start=}\FunctionTok{c}\NormalTok{(}\DecValTok{2008}\NormalTok{, }\DecValTok{12}\NormalTok{, }\DecValTok{06}\NormalTok{), }\AttributeTok{frequency=}\DecValTok{235}\NormalTok{)}
\FunctionTok{ts.plot}\NormalTok{(returns2, }\AttributeTok{main=}\StringTok{"daily return of investment index"}\NormalTok{,}
        \AttributeTok{ylab=}\StringTok{"return"}\NormalTok{, }\AttributeTok{ylim=}\FunctionTok{c}\NormalTok{(}\SpecialCharTok{{-}}\FloatTok{0.1}\NormalTok{, }\FloatTok{0.1}\NormalTok{))}
\end{Highlighting}
\end{Shaded}

\subsection{Long memory}\label{long-memory}

The GPH fractional-d estimate is about \texttt{d\ =\ 0.234}. Since this
value is between 0 and 0.5, the return series is stationary,
mean-reverting, and has long memory with finite variance.

\begin{Shaded}
\begin{Highlighting}[]
\FunctionTok{fdGPH}\NormalTok{(returns, }\AttributeTok{bandw.exp=}\FloatTok{0.5}\NormalTok{)}
\end{Highlighting}
\end{Shaded}

\subsection{ARMA order selection}\label{arma-order-selection}

The mean equation is selected by comparing AIC, BIC, and
\texttt{auto.arima}. The AIC search chooses ARMA(5,5), BIC chooses
ARMA(1,2), and \texttt{auto.arima} chooses ARIMA(4,0,1). Because the
ARMA(5,5) has the highest log-likelihood among the three reported
alternatives, it is used for the mean equation.

\begin{longtable}[]{@{}lrrr@{}}
\toprule\noalign{}
Method & p & q & Log-likelihood \\
\midrule\noalign{}
\endhead
\bottomrule\noalign{}
\endlastfoot
AIC grid search & 5 & 5 & 10851.72 \\
BIC grid search & 1 & 2 & 10835.95 \\
auto.arima & 4 & 1 & 10841.42 \\
\end{longtable}

\begin{Shaded}
\begin{Highlighting}[]
\NormalTok{AIC }\OtherTok{\textless{}{-}} \FunctionTok{matrix}\NormalTok{(}\AttributeTok{nrow=}\DecValTok{6}\NormalTok{, }\AttributeTok{ncol=}\DecValTok{6}\NormalTok{)}
\ControlFlowTok{for}\NormalTok{ (i }\ControlFlowTok{in} \DecValTok{1}\SpecialCharTok{:}\DecValTok{6}\NormalTok{) \{}
  \ControlFlowTok{for}\NormalTok{ (j }\ControlFlowTok{in} \DecValTok{1}\SpecialCharTok{:}\DecValTok{6}\NormalTok{) \{}
\NormalTok{    AIC[i,j] }\OtherTok{\textless{}{-}} \FunctionTok{AIC}\NormalTok{(}\FunctionTok{arima}\NormalTok{(returns, }\FunctionTok{c}\NormalTok{(i}\DecValTok{{-}1}\NormalTok{,}\DecValTok{0}\NormalTok{,j}\DecValTok{{-}1}\NormalTok{), }\AttributeTok{method=}\FunctionTok{c}\NormalTok{(}\StringTok{"CSS{-}ML"}\NormalTok{),}
                          \AttributeTok{optim.control=}\FunctionTok{list}\NormalTok{(}\AttributeTok{maxit=}\DecValTok{5000}\NormalTok{), }\AttributeTok{kappa=}\FloatTok{1e4}\NormalTok{), }\AttributeTok{k=}\DecValTok{2}\NormalTok{)}
\NormalTok{  \}}
\NormalTok{\}}
\ControlFlowTok{for}\NormalTok{ (i }\ControlFlowTok{in} \DecValTok{1}\SpecialCharTok{:}\DecValTok{6}\NormalTok{) \{}
  \ControlFlowTok{for}\NormalTok{ (j }\ControlFlowTok{in} \DecValTok{1}\SpecialCharTok{:}\DecValTok{6}\NormalTok{) \{}
    \ControlFlowTok{if}\NormalTok{ (AIC[i,j] }\SpecialCharTok{==} \FunctionTok{min}\NormalTok{(AIC)) \{ p }\OtherTok{\textless{}{-}}\NormalTok{ i}\DecValTok{{-}1}\NormalTok{; q }\OtherTok{\textless{}{-}}\NormalTok{ j}\DecValTok{{-}1}\NormalTok{ \}}
\NormalTok{  \}}
\NormalTok{\}}
\NormalTok{arma\_sarmaye }\OtherTok{\textless{}{-}} \FunctionTok{arima}\NormalTok{(returns, }\FunctionTok{c}\NormalTok{(p,}\DecValTok{0}\NormalTok{,q), }\AttributeTok{method=}\FunctionTok{c}\NormalTok{(}\StringTok{"CSS{-}ML"}\NormalTok{),}
                      \AttributeTok{optim.control=}\FunctionTok{list}\NormalTok{(}\AttributeTok{maxit=}\DecValTok{5000}\NormalTok{), }\AttributeTok{kappa=}\FloatTok{1e4}\NormalTok{)}
\NormalTok{arma\_sarmaye}

\NormalTok{BIC }\OtherTok{\textless{}{-}} \FunctionTok{matrix}\NormalTok{(}\AttributeTok{nrow=}\DecValTok{6}\NormalTok{, }\AttributeTok{ncol=}\DecValTok{6}\NormalTok{)}
\ControlFlowTok{for}\NormalTok{ (i }\ControlFlowTok{in} \DecValTok{1}\SpecialCharTok{:}\DecValTok{6}\NormalTok{) \{}
  \ControlFlowTok{for}\NormalTok{ (j }\ControlFlowTok{in} \DecValTok{1}\SpecialCharTok{:}\DecValTok{6}\NormalTok{) \{}
\NormalTok{    BIC[i,j] }\OtherTok{\textless{}{-}} \FunctionTok{BIC}\NormalTok{(}\FunctionTok{arima}\NormalTok{(returns, }\FunctionTok{c}\NormalTok{(i}\DecValTok{{-}1}\NormalTok{,}\DecValTok{0}\NormalTok{,j}\DecValTok{{-}1}\NormalTok{), }\AttributeTok{method=}\FunctionTok{c}\NormalTok{(}\StringTok{"CSS{-}ML"}\NormalTok{),}
                          \AttributeTok{optim.control=}\FunctionTok{list}\NormalTok{(}\AttributeTok{maxit=}\DecValTok{5000}\NormalTok{), }\AttributeTok{kappa=}\FloatTok{1e4}\NormalTok{))}
\NormalTok{  \}}
\NormalTok{\}}
\ControlFlowTok{for}\NormalTok{ (i }\ControlFlowTok{in} \DecValTok{1}\SpecialCharTok{:}\DecValTok{6}\NormalTok{) \{}
  \ControlFlowTok{for}\NormalTok{ (j }\ControlFlowTok{in} \DecValTok{1}\SpecialCharTok{:}\DecValTok{6}\NormalTok{) \{}
    \ControlFlowTok{if}\NormalTok{ (BIC[i,j] }\SpecialCharTok{==} \FunctionTok{min}\NormalTok{(BIC)) \{ p }\OtherTok{\textless{}{-}}\NormalTok{ i}\DecValTok{{-}1}\NormalTok{; q }\OtherTok{\textless{}{-}}\NormalTok{ j}\DecValTok{{-}1}\NormalTok{ \}}
\NormalTok{  \}}
\NormalTok{\}}
\NormalTok{arma\_sarmaye2 }\OtherTok{\textless{}{-}} \FunctionTok{arima}\NormalTok{(returns, }\FunctionTok{c}\NormalTok{(p,}\DecValTok{0}\NormalTok{,q), }\AttributeTok{method=}\FunctionTok{c}\NormalTok{(}\StringTok{"CSS{-}ML"}\NormalTok{),}
                       \AttributeTok{optim.control=}\FunctionTok{list}\NormalTok{(}\AttributeTok{maxit=}\DecValTok{5000}\NormalTok{), }\AttributeTok{kappa=}\FloatTok{1e4}\NormalTok{)}
\NormalTok{arma\_sarmaye2}

\FunctionTok{auto.arima}\NormalTok{(returns)}

\NormalTok{mean.equ }\OtherTok{\textless{}{-}}\NormalTok{ arma\_sarmaye}
\FunctionTok{summary}\NormalTok{(mean.equ)}
\end{Highlighting}
\end{Shaded}

\subsection{Mean-model diagnostics}\label{mean-model-diagnostics}

The ACF and PACF of the ARMA residuals are close to zero. The Box-Pierce
and Ljung-Box p-values are both approximately 0.958, so the null
hypothesis that residuals are white noise is not rejected at the 5
percent level. This supports the adequacy of the mean equation.

\begin{Shaded}
\begin{Highlighting}[]
\NormalTok{mean.residuals }\OtherTok{\textless{}{-}}\NormalTok{ mean.equ}\SpecialCharTok{$}\NormalTok{residuals}
\FunctionTok{acf}\NormalTok{(mean.residuals, }\AttributeTok{lag.max=}\DecValTok{8}\NormalTok{)}
\FunctionTok{pacf}\NormalTok{(mean.residuals, }\AttributeTok{lag.max=}\DecValTok{8}\NormalTok{)}

\FunctionTok{Box.test}\NormalTok{(mean.residuals, }\AttributeTok{lag=}\FunctionTok{log}\NormalTok{(}\DecValTok{3393}\NormalTok{), }\FunctionTok{c}\NormalTok{(}\StringTok{"Box{-}Pierce"}\NormalTok{))}
\FunctionTok{Box.test}\NormalTok{(mean.residuals, }\AttributeTok{lag=}\FunctionTok{log}\NormalTok{(}\DecValTok{3393}\NormalTok{), }\FunctionTok{c}\NormalTok{(}\StringTok{"Ljung{-}Box"}\NormalTok{))}
\end{Highlighting}
\end{Shaded}

\subsection{Distributional analysis and ARCH
effects}\label{distributional-analysis-and-arch-effects}

The return distribution is right-skewed and leptokurtic. The reported
skewness is approximately 0.730 and excess kurtosis is approximately
3.894. The Anderson-Darling normality test gives a p-value below 0.05,
so normality is rejected. Tests on squared residuals indicate remaining
volatility dependence, which motivates GARCH-type volatility modeling.

\begin{Shaded}
\begin{Highlighting}[]
\FunctionTok{skewness}\NormalTok{(returns)}
\FunctionTok{kurtosis}\NormalTok{(returns)}
\FunctionTok{hist}\NormalTok{(returns)}
\FunctionTok{ad.test}\NormalTok{(returns)}
\FunctionTok{plot}\NormalTok{(}\FunctionTok{density}\NormalTok{(returns))}
\FunctionTok{chart.QQPlot}\NormalTok{(returns, }\AttributeTok{distribution=}\StringTok{"norm"}\NormalTok{)}

\FunctionTok{acf}\NormalTok{(mean.residuals}\SpecialCharTok{\^{}}\DecValTok{2}\NormalTok{, }\AttributeTok{lag.max=}\DecValTok{30}\NormalTok{)}
\FunctionTok{pacf}\NormalTok{(mean.residuals}\SpecialCharTok{\^{}}\DecValTok{2}\NormalTok{, }\AttributeTok{lag.max=}\DecValTok{30}\NormalTok{)}
\FunctionTok{Box.test}\NormalTok{(mean.residuals}\SpecialCharTok{\^{}}\DecValTok{2}\NormalTok{, }\AttributeTok{lag=}\FunctionTok{log}\NormalTok{(}\DecValTok{3393}\NormalTok{), }\FunctionTok{c}\NormalTok{(}\StringTok{"Box{-}Pierce"}\NormalTok{))}
\FunctionTok{Box.test}\NormalTok{(mean.residuals}\SpecialCharTok{\^{}}\DecValTok{2}\NormalTok{, }\AttributeTok{lag=}\FunctionTok{log}\NormalTok{(}\DecValTok{3393}\NormalTok{), }\FunctionTok{c}\NormalTok{(}\StringTok{"Ljung{-}Box"}\NormalTok{))}
\end{Highlighting}
\end{Shaded}

\subsection{FI-GARCH modeling}\label{fi-garch-modeling}

A FI-GARCH specification is fitted with an ARMA(5,5) mean equation,
ARFIMA terms, and skewed Student-t innovations. A second model also
includes the overall market index return as an external regressor in the
mean equation.

\begin{Shaded}
\begin{Highlighting}[]
\NormalTok{garch }\OtherTok{\textless{}{-}} \FunctionTok{ugarchspec}\NormalTok{(}\AttributeTok{variance.model=}\FunctionTok{list}\NormalTok{(}\AttributeTok{model=}\StringTok{"fiGARCH"}\NormalTok{, }\AttributeTok{garchOrder=}\FunctionTok{c}\NormalTok{(}\DecValTok{1}\NormalTok{,}\DecValTok{1}\NormalTok{)),}
                     \AttributeTok{mean.model=}\FunctionTok{list}\NormalTok{(}\AttributeTok{armaOrder=}\FunctionTok{c}\NormalTok{(}\DecValTok{5}\NormalTok{,}\DecValTok{5}\NormalTok{), }\AttributeTok{include.mean=}\ConstantTok{FALSE}\NormalTok{,}
                                     \AttributeTok{archm=}\ConstantTok{FALSE}\NormalTok{, }\AttributeTok{archpow=}\DecValTok{1}\NormalTok{, }\AttributeTok{arfima=}\ConstantTok{TRUE}\NormalTok{,}
                                     \AttributeTok{archex=}\ConstantTok{FALSE}\NormalTok{),}
                     \AttributeTok{distribution.model=}\StringTok{"sstd"}\NormalTok{)}
\NormalTok{fit }\OtherTok{\textless{}{-}} \FunctionTok{ugarchfit}\NormalTok{(garch, }\AttributeTok{data=}\NormalTok{returns)}
\NormalTok{fit}
\FunctionTok{plot}\NormalTok{(fit, }\AttributeTok{which=}\StringTok{"all"}\NormalTok{)}

\NormalTok{z1 }\OtherTok{\textless{}{-}} \FunctionTok{ts}\NormalTok{(ind}\SpecialCharTok{$}\NormalTok{index)}
\NormalTok{z2 }\OtherTok{\textless{}{-}} \FunctionTok{diff}\NormalTok{(}\FunctionTok{log}\NormalTok{(z1), }\AttributeTok{lag=}\DecValTok{1}\NormalTok{)}
\NormalTok{garch2 }\OtherTok{\textless{}{-}} \FunctionTok{ugarchspec}\NormalTok{(}\AttributeTok{variance.model=}\FunctionTok{list}\NormalTok{(}\AttributeTok{model=}\StringTok{"fiGARCH"}\NormalTok{, }\AttributeTok{garchOrder=}\FunctionTok{c}\NormalTok{(}\DecValTok{1}\NormalTok{,}\DecValTok{1}\NormalTok{)),}
                      \AttributeTok{mean.model=}\FunctionTok{list}\NormalTok{(}\AttributeTok{armaOrder=}\FunctionTok{c}\NormalTok{(}\DecValTok{5}\NormalTok{,}\DecValTok{5}\NormalTok{), }\AttributeTok{include.mean=}\ConstantTok{FALSE}\NormalTok{,}
                                      \AttributeTok{archm=}\ConstantTok{FALSE}\NormalTok{, }\AttributeTok{archpow=}\DecValTok{1}\NormalTok{, }\AttributeTok{arfima=}\ConstantTok{TRUE}\NormalTok{,}
                                      \AttributeTok{archex=}\ConstantTok{TRUE}\NormalTok{,}
                                      \AttributeTok{external.regressors=}\FunctionTok{matrix}\NormalTok{(z2)),}
                      \AttributeTok{distribution.model=}\StringTok{"sstd"}\NormalTok{)}
\NormalTok{fit2 }\OtherTok{\textless{}{-}} \FunctionTok{ugarchfit}\NormalTok{(garch2, }\AttributeTok{data=}\NormalTok{returns)}
\NormalTok{fit2}
\FunctionTok{plot}\NormalTok{(fit2, }\AttributeTok{which=}\StringTok{"all"}\NormalTok{)}
\end{Highlighting}
\end{Shaded}

\subsection{Value-at-Risk}\label{value-at-risk}

The fitted volatility model is used to plot conditional risk bands. The
one-period 95 percent VaR is also computed by modified, Gaussian, and
historical methods.

\begin{longtable}[]{@{}lr@{}}
\toprule\noalign{}
Method & VaR \\
\midrule\noalign{}
\endhead
\bottomrule\noalign{}
\endlastfoot
Modified & -0.01336006 \\
Gaussian & -0.01659899 \\
Historical & -0.01464882 \\
\end{longtable}

Using the modified VaR as the benchmark, there is a 5 percent
probability of losing at least about 1.3 percent in one day when holding
the financial institutions index.

\begin{Shaded}
\begin{Highlighting}[]
\NormalTok{n }\OtherTok{\textless{}{-}} \FunctionTok{length}\NormalTok{(returns)}
\NormalTok{model.fit }\OtherTok{\textless{}{-}} \FunctionTok{ugarchfit}\NormalTok{(}\AttributeTok{spec=}\NormalTok{garch, }\AttributeTok{data=}\NormalTok{returns, }\AttributeTok{solver=}\StringTok{"solnp"}\NormalTok{)}
\NormalTok{model.fit}

\FunctionTok{qplot}\NormalTok{(}\AttributeTok{y=}\NormalTok{returns, }\AttributeTok{x=}\DecValTok{1}\SpecialCharTok{:}\NormalTok{n, }\AttributeTok{geom=}\StringTok{"point"}\NormalTok{) }\SpecialCharTok{+}
  \FunctionTok{geom\_point}\NormalTok{(}\AttributeTok{size=}\FloatTok{0.1}\NormalTok{) }\SpecialCharTok{+}
  \FunctionTok{geom\_line}\NormalTok{(}\FunctionTok{aes}\NormalTok{(}\AttributeTok{y=}\NormalTok{model.fit}\SpecialCharTok{@}\NormalTok{fit}\SpecialCharTok{$}\NormalTok{sigma}\SpecialCharTok{*}\FunctionTok{qdist}\NormalTok{(}\AttributeTok{distribution=}\StringTok{"sstd"}\NormalTok{,}
                                            \AttributeTok{shape=}\FloatTok{3.4125}\NormalTok{, }\AttributeTok{p=}\FloatTok{0.05}\NormalTok{), }\AttributeTok{x=}\DecValTok{1}\SpecialCharTok{:}\NormalTok{n),}
            \AttributeTok{colour=}\StringTok{"red"}\NormalTok{, }\AttributeTok{size=}\FloatTok{1.0}\NormalTok{) }\SpecialCharTok{+}
  \FunctionTok{geom\_line}\NormalTok{(}\FunctionTok{aes}\NormalTok{(}\AttributeTok{y=}\NormalTok{model.fit}\SpecialCharTok{@}\NormalTok{fit}\SpecialCharTok{$}\NormalTok{sigma}\SpecialCharTok{*}\FunctionTok{qdist}\NormalTok{(}\AttributeTok{distribution=}\StringTok{"sstd"}\NormalTok{,}
                                            \AttributeTok{shape=}\FloatTok{3.4125}\NormalTok{, }\AttributeTok{p=}\FloatTok{0.95}\NormalTok{), }\AttributeTok{x=}\DecValTok{1}\SpecialCharTok{:}\NormalTok{n),}
            \AttributeTok{colour=}\StringTok{"red"}\NormalTok{, }\AttributeTok{size=}\FloatTok{1.0}\NormalTok{) }\SpecialCharTok{+}
  \FunctionTok{labs}\NormalTok{(}\AttributeTok{x=}\StringTok{" "}\NormalTok{, }\AttributeTok{y=}\StringTok{"Daily Return"}\NormalTok{, }\AttributeTok{title=}\StringTok{"Value at Risk Comparison"}\NormalTok{)}

\NormalTok{sarmaye }\OtherTok{\textless{}{-}} \FunctionTok{data.frame}\NormalTok{(}\FunctionTok{cbind}\NormalTok{(}\AttributeTok{year=}\FunctionTok{substr}\NormalTok{(sarmaye[,}\DecValTok{1}\NormalTok{],}\DecValTok{1}\NormalTok{,}\DecValTok{4}\NormalTok{),}
                            \AttributeTok{month=}\FunctionTok{substr}\NormalTok{(sarmaye[,}\DecValTok{1}\NormalTok{],}\DecValTok{5}\NormalTok{,}\DecValTok{6}\NormalTok{),}
                            \AttributeTok{day=}\FunctionTok{substr}\NormalTok{(sarmaye[,}\DecValTok{1}\NormalTok{],}\DecValTok{7}\NormalTok{,}\DecValTok{8}\NormalTok{), sarmaye))}
\NormalTok{sarmaye}\SpecialCharTok{$}\NormalTok{date }\OtherTok{\textless{}{-}} \FunctionTok{as.Date}\NormalTok{(sarmaye}\SpecialCharTok{$}\NormalTok{date)}
\NormalTok{price\_xts }\OtherTok{\textless{}{-}} \FunctionTok{xts}\NormalTok{(sarmaye}\SpecialCharTok{$}\NormalTok{price, }\AttributeTok{order.by=}\NormalTok{sarmaye}\SpecialCharTok{$}\NormalTok{date)}
\NormalTok{return\_xts }\OtherTok{\textless{}{-}} \FunctionTok{diff}\NormalTok{(}\FunctionTok{log}\NormalTok{(price\_xts), }\AttributeTok{lag=}\DecValTok{1}\NormalTok{)}

\FunctionTok{mean}\NormalTok{(returns)}
\FunctionTok{sd}\NormalTok{(returns)}
\FunctionTok{skewness}\NormalTok{(returns)}
\FunctionTok{kurtosis}\NormalTok{(returns)}

\FunctionTok{VaR}\NormalTok{(}\AttributeTok{R=}\NormalTok{return\_xts, }\AttributeTok{p=}\FloatTok{0.95}\NormalTok{, }\AttributeTok{method=}\FunctionTok{c}\NormalTok{(}\StringTok{"modified"}\NormalTok{), }\AttributeTok{invert=}\ConstantTok{TRUE}\NormalTok{,}
    \AttributeTok{mu=}\FloatTok{0.001377912}\NormalTok{, }\AttributeTok{sigma=}\FloatTok{0.01093079}\NormalTok{, }\AttributeTok{m3=}\FloatTok{0.7302357}\NormalTok{, }\AttributeTok{m4=}\FloatTok{3.893968}\NormalTok{)}
\FunctionTok{VaR}\NormalTok{(}\AttributeTok{R=}\NormalTok{return\_xts, }\AttributeTok{p=}\FloatTok{0.95}\NormalTok{, }\AttributeTok{method=}\FunctionTok{c}\NormalTok{(}\StringTok{"gaussian"}\NormalTok{), }\AttributeTok{invert=}\ConstantTok{TRUE}\NormalTok{,}
    \AttributeTok{mu=}\FloatTok{0.001377912}\NormalTok{, }\AttributeTok{sigma=}\FloatTok{0.01093079}\NormalTok{, }\AttributeTok{m3=}\FloatTok{0.7302357}\NormalTok{, }\AttributeTok{m4=}\FloatTok{3.893968}\NormalTok{)}
\FunctionTok{VaR}\NormalTok{(}\AttributeTok{R=}\NormalTok{return\_xts, }\AttributeTok{p=}\FloatTok{0.95}\NormalTok{, }\AttributeTok{method=}\FunctionTok{c}\NormalTok{(}\StringTok{"historical"}\NormalTok{), }\AttributeTok{invert=}\ConstantTok{TRUE}\NormalTok{,}
    \AttributeTok{mu=}\FloatTok{0.001377912}\NormalTok{, }\AttributeTok{sigma=}\FloatTok{0.01093079}\NormalTok{, }\AttributeTok{m3=}\FloatTok{0.7302357}\NormalTok{, }\AttributeTok{m4=}\FloatTok{3.893968}\NormalTok{)}
\end{Highlighting}
\end{Shaded}

\subsection{Conclusion}\label{conclusion}

The return series is stationary but long-memory. The ARMA(5,5) mean
equation is adequate according to residual autocorrelation diagnostics,
but returns are non-normal and volatility clustering remains. FI-GARCH
modeling and VaR estimation are therefore appropriate tools for the risk
analysis.

\end{document}
